（其中包括运动学），我们也就从运动学的角度——即只讨论速度与加速度，而不讨论产生加速度的原因——来讨论三角函数。

一切研究工作都应该从最简单的情况开始。牛顿力学始于匀速直线运动，从这里有了惯性系的概念。在讨论周期运动时，我们自然从匀速圆周运动开始。

表示平面上的运动当然用平面向量为好，圆周运动则用极坐标最方便。以下我们以圆周运动的中心 $O$ 为极坐标原点（如图2-2-1），以一个指定的轴 $OA$ 为极轴，于是一个点 $P$ 决定了一个向量 $\overrightarrow{OP}$。所有的向量均分成两个分量——即轴向的与横向的。在这两个方向上的单位向量记为 $\boldsymbol\rho$ 与 $\boldsymbol\theta$，一点 $P$ 的辐角 $\varphi$ 不受任何限制：$-\infty<\varphi<+\infty$，所以 $\varphi$ 与 $\varphi+2k\pi$，$k=0,\pm1,\pm2\cdots$ 表示同一个点，这使 $\varphi$ 恰好适用于表示周期现象。而 $\varphi$ 的多值性则时常是产生问题的根源。正如大家都已习惯了的，角度用弧度表示，但是在微积分中弧度有时表示扇形面积占全圆面积的百分比（全圆的面积 $=$（半径）$^2\times\pi$弧度），有时表示角度。但是它总是无量纲量。匀速圆周运动就是角速度 $\alpha=\text{常数}$ 的情况。$\alpha$ 的单位是弧度/秒，但是我们时常以每秒绕圆心的周数 $\omega$ 度量角速度，因此其单位是周/秒，$\omega=\alpha/(2\pi)$，所以匀速圆周运动的方程是
\[
\varphi=\alpha t=2\pi\omega t. \tag{1}
\]
后面的表示法是更常用的。在一般地刻画周期运动时，$\alpha$ 称为圆频率，$\omega$ 称为频率\jiaokan{原文作“$\alpha$ 称为频率，$\omega$ 称为圆频率”。但前文已定义 $\omega=\alpha/(2\pi)$，其中 $\omega$ 表示每秒绕圆心的周数；依此定义，$\omega$ 应为普通频率，而 $\alpha$ 为圆频率（角频率）。}。它们的量纲均为 $1/T$。

\begin{figure}[htbp]
\centering
\begin{tikzpicture}[bella figure,scale=0.95,>=stealth]
  \coordinate (O) at (0,0);
  \draw (O) circle (1.65);
  \draw[->] (O)--(2.55,0) node[right] {$A$};
  \coordinate (P) at ({1.65*cos(48)},{1.65*sin(48)});
  \draw[thick,->] (O)--(P) node[pos=.55,left] {$\rho$};
  \draw[->] (P)--++({0.75*cos(138)},{0.75*sin(138)}) node[above] {$\theta$};
  \draw (0.52,0) arc[start angle=0,end angle=48,radius=0.52];
  \node at (0.72,0.25) {$\varphi$};
  \node[left] at (O) {$O$};
  \node[right] at (P) {$P$};
\end{tikzpicture}
\caption*{图2-2-1}
\end{figure}

\begin{figure}[htbp]
\centering
\begin{tikzpicture}[scale=0.86,>=stealth,
  every node/.style={font=\small,fill=none,inner sep=0pt,outer sep=1.5pt}]
  % (a): 按原图重绘；文字避开弧线与线段
  \begin{scope}[xshift=-3.25cm]
    \coordinate (Oa) at (0,0);
    \coordinate (Pa) at ({2.10*cos(35)},{2.10*sin(35)});
    \coordinate (Ppa) at ({2.10*cos(52)},{2.10*sin(52)});
    \draw[thick] (Oa) ++(-7:2.10) arc[start angle=-7,end angle=66,radius=2.10];
    \draw[->] (Oa)--(Pa);
    \draw[->] (Oa)--(Ppa);
    \draw[thick,->] (Pa)--(Ppa);
    \draw (0.64,0) arc[start angle=0,end angle=35,radius=0.64];
    \node at (0.80,0.31) {$d\varphi$};
    \node at (1.05,1.13) {$d\boldsymbol\rho$};
    \node[below left=2pt] at (Oa) {$O$};
    \node[right=3pt] at (Pa) {$P$};
    \node[above left=3pt] at (Ppa) {$P'$};
    \node at (1.03,-0.50) {(a)};
  \end{scope}
  % (b)
  \begin{scope}[xshift=3.25cm]
    \coordinate (Ob) at (0,0);
    \coordinate (Pb) at ({2.10*cos(35)},{2.10*sin(35)});
    \coordinate (Ppb) at ({2.10*cos(49)},{2.10*sin(49)});
    \draw[thick] (Ob) ++(-7:2.10) arc[start angle=-7,end angle=66,radius=2.10];
    \draw[->] (Ob)--(Pb);
    \draw[->] (Ob)--(Ppb);
    % 两个相邻切向单位向量
    \coordinate (T) at ($(Pb)+({1.32*cos(125)},{1.32*sin(125)})$);
    \coordinate (Tp) at ($(Ppb)+({1.32*cos(139)},{1.32*sin(139)})$);
    \draw[->] (Pb)--(T);
    \draw[->] (Ppb)--(Tp);
    \draw[thick,->] (T)--(Tp);
    % 原图中的辅助点 R 与平移后的短向量
    \coordinate (R) at ($(Pb)+(Tp)-(T)$);
    \draw[thick,->] (Pb)--(R);
    \draw (0.64,0) arc[start angle=0,end angle=35,radius=0.64];
    \node at (0.80,0.31) {$d\varphi$};
    \node[below left=2pt] at (Ob) {$O$};
    \node[right=3pt] at (Pb) {$P$};
    \node[left=3pt] at (Ppb) {$P'$};
    \node[above right=5pt] at (T) {$T$};
    \node[above left=7pt] at (Tp) {$T'$};
    \node[below left=5pt] at (R) {$R$};
    \node[above=9pt] at ($(T)!0.50!(Tp)$) {$d\boldsymbol\theta$};
    \node[right=7pt] at ($(Pb)!0.62!(T)$) {$\boldsymbol\theta$};
    \node at (1.03,-0.50) {(b)};
  \end{scope}
\end{tikzpicture}
\caption*{图2-2-2}
\end{figure}

下面我们在极坐标中计算运动学的主要变量：速度和加速度，因为我们要在这个基础上建立三角函数理论，所以请读者注意，这里完全不许使用三角函数知识。若 $P$ 点之极坐标为 $(r,\varphi)$，注意到随着时间的变化，$\boldsymbol\rho$ 与 $\boldsymbol\theta$ 大小虽然都不变（$|\boldsymbol\rho|=|\boldsymbol\theta|=1$），但是其方向是改变的，所以，由 $\overrightarrow{OP}=r\boldsymbol\rho$，对 $t$ 求导后可以得到
\[
\boldsymbol v=\frac{d}{dt}\overrightarrow{OP}=\frac{dr}{dt}\boldsymbol\rho+r\frac{d\boldsymbol\rho}{dt}
=\dot r\boldsymbol\rho+r\dot{\boldsymbol\rho}. \tag{2}
\]
这里我们按牛顿的作法，用上加一点表示对时间 $t$ 求导。

现在来计算 $\dot{\boldsymbol\rho}=d\boldsymbol\rho/dt$。如图2-2-2的(a)作单位圆，若 $\boldsymbol\rho$ 记为 $\overrightarrow{OP}$，经过时间 $dt$ 后，若 $P$ 沿单位圆变成 $P'$，则 $\boldsymbol\rho$ 有一个增量 $d\boldsymbol\rho$ 即 $\overrightarrow{PP'}$，它对于圆心张一个角 $d\varphi$，因此按弧度计，利用 $|\overrightarrow{OP}|=1$，即知 $d\boldsymbol\rho$ 之大小即 $d\varphi$，而其方向则是切线方向 $\overrightarrow{PP'}$，即 $\boldsymbol\theta$ 方向。总之
\[
\dot{\boldsymbol\rho}=\frac{d\boldsymbol\rho}{dt}=\frac{d\varphi}{dt}\boldsymbol\theta=\dot\varphi\boldsymbol\theta. \tag{3}
\]
$\dot{\boldsymbol\rho}$ 的方向恰好与 $\boldsymbol\rho$ 正交。

再看加速度 $\boldsymbol a$。由(2)
\[
\begin{aligned}
\boldsymbol a=\dot{\boldsymbol v}
&=\ddot r\boldsymbol\rho+2\dot r\dot{\boldsymbol\rho}+r\frac{d\dot{\boldsymbol\rho}}{dt}\\
&=\ddot r\boldsymbol\rho+2\dot r\dot\varphi\boldsymbol\theta+r(\ddot\varphi\boldsymbol\theta+\dot\varphi\dot{\boldsymbol\theta}).
\end{aligned}\tag{4}
\]
所以现在又要计算 $\dot{\boldsymbol\theta}$，再看（图2-2-2,b）的单位圆。若在无穷小时间 $dt$ 内 $\overrightarrow{OP}$ 变成了 $\overrightarrow{OP'}$，$\boldsymbol\theta$ 由 $P$ 点处切线上的单位向量 $\overrightarrow{PT}$ 变成 $P'$ 处切线上的单位向量 $\overrightarrow{P'T'}$，而 $d\boldsymbol\theta$ 就是向量 $\overrightarrow{TT'}$，将 $\overrightarrow{P'T'}$ 平移为 $\overrightarrow{PR}$，则 $\overrightarrow{TT'}=\overrightarrow{PR}$，而由于 $PT\perp OP$，$P'T'\perp OP'$，而且 $|OP|=|PT|=1$，那么，$\triangle POP'\cong\triangle TPR$。这样，当 $d\varphi$ 为无穷小量时，$TT'\perp PT$ 因而指向圆心 $O$ 而且方向相反，这样
\[
\dot{\boldsymbol\theta}=\frac{d\boldsymbol\theta}{dt}=-\frac{d\varphi}{dt}\boldsymbol\rho=-\dot\varphi\boldsymbol\rho. \tag{5}
\]
代入(4)式，即得
\[
\begin{aligned}
\boldsymbol a
&=\ddot r\boldsymbol\rho+(2\dot r\dot\varphi+r\ddot\varphi)\boldsymbol\theta-r\dot\varphi^2\boldsymbol\rho\\
&=(\ddot r-r\dot\varphi^2)\boldsymbol\rho+(2\dot r\dot\varphi+r\ddot\varphi)\boldsymbol\theta.
\end{aligned}\tag{6}
\]
在这个计算过程中，我们时常容许一些误差。但是读者会看到，若以 $d\varphi$ 表示一阶无穷小量，则我们略去的都是关于 $d\varphi$ 的高阶无穷小量。在第三章中，我们将指出，略去高阶无穷小量是微分学的基本思想，而且全章就是为了说明这个思想怎样由它最初的形态，发展到它的现代的形式，构成整个微分学。所以现在我们就不去讲它了。

以上讲的适用于一切运动。如果把它应用于圆周运动，而以该圆周的中心为极坐标原点，将得到 $r=\text{常数}$，从而 $\dot r=0,\ddot r=0$（但是 $\boldsymbol\rho$ 因为方向会变所以不是常向量），而(2)式变成
\[
\boldsymbol v=r\dot\varphi\boldsymbol\theta. \tag{2'}
\]
这就是说，圆周运动只有切向速度，大小为 $|r\dot\varphi|$ 即我们所熟悉的线速度 $=$ 半径 $\times$ 角速度。(6)则成为
\[
\boldsymbol a=-r\dot\varphi^2\boldsymbol\rho+r\ddot\varphi\boldsymbol\theta. \tag{6'}
\]
它由两部分组成：一是径向加速度 $-r\dot\varphi^2\boldsymbol\rho$，亦即我们常说的向心加速度，另一部分是切向的线加速度 $=$ 半径 $\times$ 角加速度。

再看匀速圆周运动，这时 $\dot\varphi=\alpha=\text{常数}$，于是 $\ddot\varphi=0$ 而由(6')有
\[
\frac{d^2}{dt^2}\overrightarrow{OP}=-\alpha^2\overrightarrow{OP}. \tag{6''}
\]
这是一个二阶常微分方程。下面我们所要作的事就是从这样的方程来定义三角函数。这里我们还要提醒一下：推导出(6'')时，我们只利用了一些运动学的概念，即时间、空间、速度与加速度，而完全没有利用动力学的概念如力、功、能量等等。下面我们还会把这些动力学概念用于方程(6'')的研究。

\textbf{2. 指数函数与三角函数}\quad 平面向量都可以用复数来表示。所以现在我们用 $z=z(t)$ 来表示 $\overrightarrow{OP}$，而方程(6'')就成了
\[
\ddot z(t)+\alpha^2z(t)=0. \tag{7}
\]
如果分开实虚部：$z(t)=x(t)+\sqrt{-1}\,y(t)$ 就有
\[
\ddot x(t)+\alpha^2x(t)=0, \tag{8$_1$}
\]
\[
\ddot y(t)+\alpha^2y(t)=0. \tag{8$_2$}
\]
我们通常用 $i$ 表示 $\sqrt{-1}$，但有的物理学家反对，他们认为 $i$ 通常表示例如电流，所以主张用 $j$ 代替之，但还有物理学家反唇相讥说 $j$ 有时也表示某个物理量，例如电流密度。所以大家都同意，在有可能引起疑义时，就用 $\sqrt{-1}$ 表示虚数单位。不论如何，这样一来，复数进入了。这时我们只要注意一点：把 $\sqrt{-1}$ 当作常数处理，因而可以进入或移出微分与积分号外。这不是一个只图一时方便的约定，而有很深的道理。在讨论线性空间时，我们应该讲明，“系数”在哪一个域（域就是可以进行四则运算的地方）中。是实数域还是复数域，线性空间的基本定理都不会改变（但是欧氏空间相应的对象叫埃尔米特空间）。所以，几何复数就意味着现在最好把复数域作为“基域”，基域中的元素就是常数。(7)或者(8)是二阶常微分方程。即令在复数域中其柯西问题之解的存在与唯一性定理仍成立，因为一个复的常微分方程如(7)，在分开实虚部后就成为一组两个实的常微分方程如(8$_1$),(8$_2$)。不过这时的初始条件应该是给定初始位置与初速，即给出
\[
z(0)=A,\qquad \dot z(0)=B. \tag{9}
\]
不过 $A$ 与 $B$ 现在都是复数。

为了求解方程(7)，我们不妨与 §1 中求解指数方程 $dE/dt=\lambda E$ 来比较，在那里我们发现了其解为 $E(t)=Ce^{\lambda t}$，就是说指数函数与其导数为“同类项”，既然如此，与其高阶导数自然也是“同类项”。所以我们不妨试一下看(7)是否有一解可以写为 $z(t)=Ce^{\lambda t}$？不过 $\lambda$ 与 $C$ 现在都应该是复数。好在上面我们已指出，哪怕复数也可以当成常数处理。故以 $Ce^{\lambda t}$ 代入(7)，即得
\[
C(\lambda^2+\alpha^2)=0.
\]
因此 $C$ 仍为未定常数，而 $\lambda=\pm\sqrt{-1}\,\alpha$，因此得到(7)的两个解
\[
z_1(t)=e^{\sqrt{-1}\alpha t},\qquad z_2(t)=e^{-\sqrt{-1}\alpha t}=\overline{z_1},
\]
注意 $z_1$ 与 $z_2$ 并不线性相关，即使在复数域中也不是。所以这就给出了(7)的基本解系。而(7)之一切解均可表示为
\[
z(t)=C_1z_1(t)+C_2z_2(t).
\]
这里 $C_1$ 与 $C_2$ 是任意复数。只要适当选取 $C_1,C_2$ 就能满足任意的初始条件(9)。

分开 $z_1(t)$ 之实部虚部并引入新记号，即得
\[
z_1(t)=C(t)+\sqrt{-1}S(t).
\]
注意到 $z_1(0)=1=1+\sqrt{-1}\cdot0$，$\dot z_1(0)=\sqrt{-1}\alpha=0+\sqrt{-1}\cdot\alpha$，即知
\[
C(0)=1,\qquad \dot C(0)=0;
\]
\[
S(0)=0,\qquad \dot S(0)=\alpha.
\]
因此关于 $C(t)$ 与 $S(t)$ 立即有
\[
\left\{\begin{aligned}
\ddot C+\alpha^2C&=0,\\
C(0)&=1,\quad \dot C(0)=0,
\end{aligned}\right. \tag{10$_1$}
\]
\[
\left\{\begin{aligned}
\ddot S+\alpha^2S&=0,\\
S(0)&=0,\quad \dot S(0)=\alpha.
\end{aligned}\right. \tag{10$_2$}
\]
我们立刻就会看到
\[
C(t)=\cos\alpha t,\qquad S(t)=\sin\alpha t.
\]
这是因为：由唯一性定理，自然有
\[
e^{\sqrt{-1}\alpha t}=\cos\alpha t+\sqrt{-1}\sin\alpha t. \tag{11}
\]
(11)式是极重要的欧拉公式，而且公认是最漂亮的数学公式之一。我们还想多说一些话。$z(t)=\cos\alpha t+\sqrt{-1}\sin\alpha t$ 为什么成立？我们再细细分析一下。这个式子之成立是我们实际验证的结果，验证中的关键步骤又是证明 $d\sin\theta/d\theta=\cos\theta$。证明这个求导公式当然不会用到(11)式，因此没有循环论证之嫌。但是它用到了微积分中另一个重要的极限式
\[
\lim_{\theta\to0}\frac{\sin\theta}{\theta}=1. \tag{12}
\]
这个极限式与另一个极限式 $\lim_{x\to0}(e^x-1)/x=1$ 同时被称为两个最重要的极限式。既然已经有了(11)式，当然就会问，这两个极限式的关系是什么？我们先把话说在前面：(12)式其实与后一极限是同一回事。这一点我们慢慢道来。且先问一下，(12)式在微积分教本中是怎样证的？它的基础是，当 $\theta$ 充分小时（不妨设它为正）
\[
\sin\theta<\theta<\tan\theta. \tag{13}
\]
这里 $\theta$ 要按弧度计，因为若把弧度按弧长来解释，这个不等式就是
\[
\text{弦长}<\text{弧长}<\text{切线长}.
\]
但是如果再追究一步，何以知道这个不等式成立就会发现它很不好证。因此在微积分教本中至少在这里是用扇形面积来解释弧度的。把问题归结为面积有很深的含意，我们暂时把它放在一边，而承认(13)式，用 $\sin\theta$（这是正数）通除(13)式，又有
\[
1<\frac{\theta}{\sin\theta}<\frac1{\cos\theta}.
\]
于是只要证最后一项极限（$\theta\to0$ 时）为1即可，用倍角公式
\[
|\cos\theta-1|=\left|\cos\left(2\cdot\frac\theta2\right)-1\right|
=\left|1-2\sin^2\frac\theta2-1\right|
=2\sin^2\frac\theta2<2\left(\frac\theta2\right)^2
\]
（利用 $\sin\frac\theta2<\frac\theta2$），故当 $\theta\to0$ 时，$\cos\theta\to1$，而基本的极限式(12)得证。

倍角公式的来源是和角公式，例如
\[
\cos(\theta_1+\theta_2)=\cos\theta_1\cos\theta_2-\sin\theta_1\sin\theta_2.
\]
它的证明当
\[
0<\theta_1,\theta_2<\frac\pi2,\qquad \theta_1+\theta_2<\frac\pi2 \tag{14}
\]
时确实是十分容易的。因为如图2-2-3
\[
\begin{aligned}
\cos(\theta_1+\theta_2)&=\frac{OC}{OA}=OC\\
&=OD-CD=OB\cos\theta_2-AB\sin\angle BAE\\
&=OA\cos\theta_1\cos\theta_2-OA\sin\theta_2\sin\theta_1\\
&=\cos\theta_1\cos\theta_2-\sin\theta_1\sin\theta_2\qquad(OA=1).
\end{aligned}
\]
这里只要注意到 $\angle BAE=\theta_2$ 即可，但是即使在中学里学三角学时，许多老师以及细心的同学都知道，如果条件(14)不成立，则不但是证明和角公式，甚至画一个适当的图也非易事。这一切告诉我们，经典的中学教本讲三角函数的办法，不仅谈不上简洁（更不谈优美了），而且相当烦琐，使人难得要领，更不说发展其本质了。所以上面讲的一切，虽然是对的，我们宁可把它放在一边，进而用其它方法证明 $C(t),S(t)$ 确实是余弦、正弦了。目前，尽管我们心里明白结果是这样，却不准用这个结果，这样我们需要重新解释
\[
e^{\sqrt{-1}t}=C(t)+\sqrt{-1}S(t) \tag{15}
\]
（为简单起见，我们取 $\alpha=1$）。

\begin{figure}[htbp]
\centering
\begin{tikzpicture}[bella figure,scale=0.85]
  \coordinate (O) at (0,0); \coordinate (A) at (1.6,2.5); \coordinate (B) at (2.35,1.65);
  \coordinate (C) at (1.6,0); \coordinate (D) at (2.35,0); \coordinate (E) at (1.6,1.65);
  \draw (O)--(A)--(B)--(O); \draw (A)--(C); \draw[dashed] (B)--(D); \draw (E)--(B);
  \node[left] at (O) {$O$}; \node[above] at (A) {$A$}; \node[right] at (B) {$B$};
  \node[below] at (C) {$C$}; \node[below] at (D) {$D$}; \node[left] at (E) {$E$};
  \node at (0.62,0.55) {$\theta_2$}; \node at (0.82,1.28) {$\theta_1$}; \node[left] at (0.72,1.42) {$1$};
\end{tikzpicture}
\caption*{图2-2-3}
\end{figure}

现在我们回到匀速圆周运动，$z(t)=e^{\sqrt{-1}t}$ 就是一个在单位圆周上以匀速 $\alpha(=1)$ 运动的动点，而 $C(t),S(t)$ 正是 $z(t)$（亦即向量 $\overrightarrow{OP}$）在坐标轴上的投影。$\overrightarrow{OP}$ 与 $x$ 轴的交角 $\theta=\arg(z)$ 是 $z$ 的辐角。这样一来我们就采用了在中学时就应该已经讲过的一般角 $\theta$（不一定限于 $0<\theta<\pi/2$）的正弦、余弦定义了：
\[
C(t)=\cos t,\qquad S(t)=\sin t.
\]
这样做的好处就是我们不必再受(14)式那样的限制。当然这意味着如和角公式这样简单的东西都得重新证。这样也好。我们可以系统地用(15)式以及 $e^{\sqrt{-1}t}$ 满足方程(7)和初值条件(10)来得出有关正余弦函数的全部性质，以此强调我们是走一条与中学数学不同的路。我们现在讲法与中学生的讲法最重要的区别何在呢？把一个向量分解成 $x$ 与 $y$ 轴两个方向的分向量之和，这并没有什么出奇之处。可是当我们用运动学的思想来看待它时，就发现匀速圆周运动有向心加速度这个简单事实就在成了一个常微分方程(7)，而且发现匀速圆周运动可以用 $e^{\sqrt{-1}t}$ 来表示。这样向量之分解为分向量就变成了欧拉公式(11)。于是，这个公认的数学中最漂亮的公式从运动学的观点看来，原来简单如此！因为想法不同，所以还保留记号 $C(t)$ 与 $S(t)$ 一直到最后再改回习惯的 $\cos t,\sin t$。

圆具有近乎完美无缺的对称性（直线当然也如此，不过直线可以说是半径为无穷大的圆）。对称性是整个数学的最基本的基础之一。我们当然现在还不能讲得太玄，但就我们“看得见”的圆的对称性，立即可得出：
\[
|z(t)|^2=C^2(t)+S^2(t)=1.
\]
当然您可以说这就是勾股定理或者什么别的东西。要紧的是，不论辐角多少，$z(t)$ 之长总是不变的，这当然是对称性。其次看“特别角”的正余弦之值。因为是匀速 $\alpha=1$ 旋转，所以每 $2\pi$ 个单位时间回转一周（所以 $2\pi$ 即周期），又因为旋转是匀速的，所以在 $1/4$ 个周期中必转过 $1/4$ 个单位圆。这样在 $t=0,\pi/2,\pi,3\pi/2,2\pi$ 时 $z(t)$ 位于 $A,B,C,D,A$ 诸点。由对称性和它在 $x$ 或 $y$ 轴上之投影或者为0，或者为1，总之在这些特别角上，依次有
\[
e^{\sqrt{-1}\alpha t}=1,\sqrt{-1},-1,-\sqrt{-1},1.
\]
再看圆对 $x$ 轴、$y$ 轴（以及原点）之对称性，依次得到 $P$ 之对称点 $Q,R,S$（图2-2-4），其辐角分别为 $\pi-t,\pi+t,-t$；由图即见
\[
\begin{aligned}
C(\pi-t)&=-C(t),&\quad C(\pi+t)&=-C(t),&\quad C(-t)&=C(t),\\
S(\pi-t)&=S(t),& S(\pi+t)&=-S(t),& S(-t)&=-S(t).
\end{aligned}\tag{16}
\]
(16)每行中最后的公式即 $C(t)$ 与 $S(t)$ 分别是偶函数与奇函数。最后再加上周期性
\[
e^{\sqrt{-1}(t+2k\pi)}=e^{\sqrt{-1}t},\qquad C(t+2k\pi)=C(t),\qquad S(t+2k\pi)=S(t). \tag{17}
\]
以上诸公式我们列为一组。

\begin{figure}[htbp]
\centering
\begin{tikzpicture}[bella figure,scale=0.9,>=stealth]
  \draw[->] (-2.1,0)--(2.3,0) node[right] {$x$};
  \draw[->] (0,-2.1)--(0,2.3) node[above] {$y$};
  \draw (0,0) circle (1.65);
  \coordinate (P) at ({1.65*cos(28)},{1.65*sin(28)});
  \coordinate (Q) at ({1.65*cos(152)},{1.65*sin(152)});
  \coordinate (R) at ({1.65*cos(208)},{1.65*sin(208)});
  \coordinate (S) at ({1.65*cos(332)},{1.65*sin(332)});
  \draw (0,0)--(P); \draw (0,0)--(Q); \draw (0,0)--(R); \draw (0,0)--(S);
  \node[right] at (P) {$P$}; \node[left] at (Q) {$Q$}; \node[left] at (R) {$R$}; \node[right] at (S) {$S$};
  \node[below left] at (0,0) {$O$};
  \node[right] at (1.65,0) {$A$}; \node[above] at (0,1.65) {$B$};
  \node[left] at (-1.65,0) {$C$}; \node[below] at (0,-1.65) {$D$};
  \node at (0.52,0.18) {$t$}; \node at (0.34,0.62) {$\pi-t$};
\end{tikzpicture}
\caption*{图2-2-4}
\end{figure}

第二组性质与“相位”(phase)有关。什么叫相位，并没有明确的答案，但是此词是从天文学中借来的。例如月亮的圆缺，阴历初一，天上完全没有月亮，是为朔日，然后有上弦、新月如芽，直到望日（阴历十五）满月，然后是下弦月。终于月亮完全消失，又到了下一个朔望月，这就是月相。某一个周期内月亮处于什么“阶段”——阶段英文也可叫 phase。但是为了讲 $C(t),S(t)$，还是用月亮在24小时内的升降——这是由地球自转造成的——来解释更好。如果把图2-2-4中的圆看作地球的剖面，我们说 $C(t),S(t)$ 是在坐标轴上的投影，可是坐标轴的选取有相当大的任意性：$A$ 点的水平轴（指向地平线）是 $y$ 轴，而铅直轴（指向天顶）则是 $x$ 轴，换到 $B$ 点，$B$ 的水平轴恰好是 $A$ 的铅直轴，而 $B$ 的铅直轴则是 $A$ 的水平轴。因此 $A,B$ 两点的 $C(t),S(t)$（各加下标 $A,B$ 以示区别）是互换的。如果有人在 $A$ 点，他的时间用 $t$ 表示，于是在 $t=0$（傍晚6时）月亮出来了；月出于东山之上，徘徊于斗牛之间，他看见的月亮高度 $S(t)$ 由 $S(0)=0$ 逐步变成 $S(t)>0$。直到 $t=\pi/2$（半夜12时）月正中天，$S(\pi/2)=1$。如果 $B$ 处还有一个人，他的时间是 $\tau$，则当 $t=\pi/2$ 时，$\tau=0$，于是 $B$ 看见月亮升起到地平线上，即 $S(0)=0$。$\tau=\pi/2$ 时，对 $B$ 说来恰好是半夜，而对 $A$ 说来已是 $t=\pi$，到了天亮了，月落乌啼霜满天，$S(t)$ 又变成了0。总之 $\tau=t-\pi/2$，这时我们说 $t$ 的相位“提前”$\pi/2$，或者由 $t=\tau+\pi/2$，我们也可以说 $\tau$ 有了“滞后”$\pi/2$。由图2-2-4看到 $S_A(t)=C_B(t),C_A(t)=-S_B(t)$，所以，以 $A$ 为准（略去下标）
\[
S(t)=C\left(t-\frac\pi2\right),\qquad C(t)=-S\left(t-\frac\pi2\right), \tag{18$_1$}
\]
或者以 $B$ 为准，同样略去下标，又得到
\[
S\left(t+\frac\pi2\right)=C(t),\qquad C\left(t+\frac\pi2\right)=-S(t). \tag{18$_2$}
\]
由此可见，$C(t)$ 和 $S(t)$ 其实是一回事，无非有一个相位差 $\pi/2$ 而已。至于是“提前”还是“滞后”，就看对哪一个函数而言。例如对函数 $y=f(t)$，$y=f(t-a),a>0$ 的图像就在其右方相距 $a$ 处，所以是 $f(t)$ 的相位“提前”，同样 $y=f(t+a),a>0$，就比 $y=f(t)$ 的相位“滞后”。名词并不重要，重要的是 $\pi/2$ 与 $-\pi/2$ 决不可混淆。

把(18$_1$)与(16)的最后一个式子（奇偶性）合起来，就得到中学生都熟悉的
\[
\sin\left(\frac\pi2-t\right)=\cos t,\qquad \cos\left(\frac\pi2-t\right)=\sin t.
\]
而且在直角三角形中，大家都知道这就是两个锐角互余。现在我们不限于直角三角形，而且看见一大堆公式，其实都是相位差的问题。例如
\[
\cos(\pi-t)=\cos\left(\frac\pi2+\frac\pi2-t\right)
=-\sin\left(\frac\pi2-t\right)=-\cos t,
\]
等等，都有多种方法“证明”或“变形”，可是真正重要的是相位问题。

可是相位问题的重要性还不止于此。例如对于(3)式 $\dot{\boldsymbol\rho}=\sqrt{-1}\dot\varphi\boldsymbol\rho$，如果令 $\boldsymbol\rho$ 为 $e^{\sqrt{-1}\alpha t}$，其中 $\alpha>0$，$\varphi(t)=\sqrt{-1}\alpha t$，$d\varphi=\sqrt{-1}\alpha dt$，而用复数表示，(3)就成为
\[
\frac d{dt}e^{\sqrt{-1}\alpha t}=\sqrt{-1}\alpha e^{\sqrt{-1}\alpha t}
=\alpha e^{\sqrt{-1}(\alpha t+\pi/2)}. \tag{19}
\]
所以，$z=e^{\sqrt{-1}\alpha t}$ 求导后，一方面“振幅”放大 $|\alpha|$ 倍，另一方面相位滞后 $\pi/2$（如果 $\alpha<0$，就应把它写为 $e^{-\sqrt{-1}|\alpha|t}$，这时仍会得到“振幅”（即模 $|z|$）放大 $|\alpha|$ 倍，但相位会滞后 $3\pi/2$）。与此类似，若对 $z$ 乘以某个复数 $ae^{\sqrt{-1}\varphi}$，而 $|a|>0$，则不但有 $z$ 的振幅会放大 $a$ 倍，而且也有相位滞后 $\varphi$。这些现象在物理学中十分重要。

三角函数再一类性质与加法定理有关。现在我们已经看见了，三角函数通过欧拉公式其实归结为指数函数。而指数函数最引人注目的性质是加法定理，那么加法定理在三角函数上有什么体现呢？看 $e^{\sqrt{-1}\varphi_1}$ 与 $e^{\sqrt{-1}\varphi_2}$，一方面由欧拉公式
\[
e^{\sqrt{-1}(\varphi_1+\varphi_2)}=C(\varphi_1+\varphi_2)+\sqrt{-1}S(\varphi_1+\varphi_2).
\]
另一方面由指数函数的加法定理有
\[
e^{\sqrt{-1}(\varphi_1+\varphi_2)}=e^{\sqrt{-1}\varphi_1}\cdot e^{\sqrt{-1}\varphi_2}
=[C(\varphi_1)+\sqrt{-1}S(\varphi_1)][C(\varphi_2)+\sqrt{-1}S(\varphi_2)].
\]
所以这两个式子双方必相等。分开其实、虚部立即得到三角函数的加法定理
\[
C(\varphi_1+\varphi_2)=C(\varphi_1)C(\varphi_2)-S(\varphi_1)S(\varphi_2), \tag{20$_1$}
\]
\[
S(\varphi_1+\varphi_2)=S(\varphi_1)C(\varphi_2)+C(\varphi_1)S(\varphi_2). \tag{20$_2$}
\]
有了加法定理自然也就有了倍角公式、半角公式等等。

现在我们看到一个现象：同一个结果可以用多种不同方法证明。条条大路通罗马，使人简直不知所从了。例如，$dC(t)/dt=-S(t),dS(t)/dt=C(t)$，就有多种方法证明：用指数函数的求导公式加上欧拉公式：
\[
\frac d{dt}e^{\sqrt{-1}t}=\sqrt{-1}e^{\sqrt{-1}t}=-S(t)+\sqrt{-1}C(t),
\]
\[
\frac d{dt}e^{\sqrt{-1}t}=\frac d{dt}[C(t)+\sqrt{-1}S(t)]
=\frac d{dt}C(t)+\sqrt{-1}\frac d{dt}S(t).
\]
比较双方实虚部也可以得到同样结果。

我们还可以用相位滞后加上欧拉公式来证。因为上面我们已经得到了指数函数求导后相位滞后 $\pi/2$：
\[
\frac d{dt}e^{\sqrt{-1}t}=e^{\sqrt{-1}(t+\pi/2)}
=C\left(t+\frac\pi2\right)+\sqrt{-1}S\left(t+\frac\pi2\right).
\]
另一方面又有三角函数的相位滞后公式(18)：$C(t+\pi/2)=-S(t),S(t+\pi/2)=C(t)$，代入上式又一次得到 $dC(t)/dt=-S(t),dS(t)/dt=C(t)$。

同一个结果可以用多种不同方法得到，这当然不是偶然的事。这后面是否还隐藏着更深刻的东西呢？是的，本书第五章全章就是讨论这里的问题。

总结以上的讨论就可以看到，我们的出发点就是指数函数是一个微分方程的解。从运动学的角度来看待这些问题之所以取得成效，也是因为匀速圆周运动的向心加速度公式就是 $z(t)$ 所应满足的微分方程。所以我们是通过 $e^{\sqrt{-1}t}$ 所满足的微分方程得出了三角函数的性质而完全没有利用直角三角形的边角关系。但是 $C(t),S(t)$ 也直接地是常微分方程柯西问题(10)的解（我们为简单计是限于 $\alpha=1$ 的情况），能否直接由它（仍令 $\alpha=1$）得出 $C(t),S(t)$ 的性质呢？这里又遇到了上面说的“条条道路通罗马”的状况：上面讲的许多性质都可以这样得出。但为此我们先把(10)化成一个方程组讨论起来更容易。以(10$_2$)为例，令 $\dot S(t)=\overline C(t)$，读者会感到，这不是利用到 $dS(t)/dt=d\sin t/dt=\cos t=C(t)$ 吗？何必又设它是 $\overline C(t)$ 呢？注意，前面证明 $dS/dt=C(t)$ 用的是相位和指数函数的微分方程，我们现在则想用(10)，所以暂时还不知道 $\dot S(t)$ 是什么，只好写为 $\overline C(t)$，于是有
\[
\dot{\overline C}(t)=\ddot S(t)=-S(t),
\]
\[
\ddot{\overline C}(t)=-\dot S(t)=-\overline C(t).
\]
利用它们一方面得出 $\ddot{\overline C}(t)+\overline C(t)=0$，另一方面，令 $t=0$ 又得 $\overline C(0)=1,\dot{\overline C}(0)=0$，总之 $\overline C(t)$ 是以下柯西问题之解：
\[
\ddot{\overline C}(t)+\overline C(t)=0,\qquad \overline C(0)=1,\quad \dot{\overline C}(0)=0.
\]
但这就是(10$_1$)（令 $\alpha=1$），于是由唯一性定理知 $\overline C(t)=C(t)$，这也算是 $d\sin t/dt=\cos t$ 的“另证”吧！

总之我们有了一个常微分方程组
\[
\dot C(t)=-S(t),\qquad \dot S(t)=C(t),
\]
\[
C(0)=1,\qquad S(0)=0. \tag{21}
\]
有了它，许多性质的证明都十分容易了，例如 $C^2(t)+S^2(t)=1$ 的证明如下：
\[
\frac d{dt}[C^2(t)+S^2(t)]=2[C(t)\dot C(t)+S(t)\dot S(t)]
=2[-C(t)S(t)+C(t)S(t)]=0.
\]
所以，$C^2(t)+S^2(t)$ 之值不随 $t$ 变化，因此
\[
C^2(t)+S^2(t)=C^2(0)+S^2(0)=1.
\]
又如加法定理的证明，注意到
\[
\begin{aligned}
\frac d{dt}[S(t)C(a-t)+C(t)S(a-t)]
&=C(t)C(a-t)+S(t)S(a-t)\\
&\quad-S(t)S(a-t)-C(t)C(a-t)=0.
\end{aligned}
\]
所以 $S(t)C(a-t)+C(t)S(a-t)=S(0)C(a-0)+C(0)S(a-0)=S(a)$。令 $a=x+y,t=x$，即得
\[
S(x+y)=S(x)C(y)+C(x)S(y).
\]
这样“证”下去，不免令人感到烦琐而无新意，但是下面我们就立刻会得到引人注目的结果了。上一节我们给出了指数函数 $e^t$ 的幂级数展开式
\[
e^t=\sum_{n=0}^{\infty}\frac{t^n}{n!}.
\]
如果把 $t$ 改成 $\sqrt{-1}t$，然后在其右方开实部与虚部，将有
\[
\begin{aligned}
e^{\sqrt{-1}t}&=C(t)+\sqrt{-1}S(t)\\
&=\left(1-\frac{t^2}{2!}+\frac{t^4}{4!}-\cdots\right)
+\sqrt{-1}\left(t-\frac{t^3}{3!}+\frac{t^5}{5!}-\cdots\right).
\end{aligned}
\]
所以
\[
S(t)=t-\frac{t^3}{3!}+\frac{t^5}{5!}-\cdots
=\sum_{n=0}^{\infty}(-1)^n\frac{t^{2n+1}}{(2n+1)!}, \tag{22$_1$}
\]
\[
C(t)=1-\frac{t^2}{2!}+\frac{t^4}{4!}-\cdots
=\sum_{n=0}^{\infty}(-1)^n\frac{t^{2n}}{(2n)!}. \tag{22$_2$}
\]
这两个幂级数对一切 $t$（不论实数还是复数）都是对的，而且可以逐项微分、逐项积分。我们也不妨直接把(22$_1$)与(22$_2$)代入方程组(21)及相应的初始条件就知道这两个级数确实适合它们。

现在要回答遗留下来的一个问题：不利用直角三角形边角关系如何证明 $\lim_{x\to0}\sin x/x=1$？一方面可以说它太容易，不足道了；用级数(22$_1$)立即可知，但是这个证明显然“不好”，因为它没有给我们任何新的启发，所以我们来看另一条“道路”。微积分中另一个极重要的极限式 $\lim_{t\to0}(e^t-1)/t=1$ 是怎么证的？当时我们利用 $e^0=1$，可知
\[
\frac{e^t-1}{t}=\frac{e^{0+t}-e^0}{t}=\left.\frac{\Delta e^x}{\Delta x}\right|_{x=0,\Delta x=t}. \tag{23}
\]
于是这个极限就是 $de^x/dx|_{x=0}=e^x|_{x=0}=1$。注意，我们有意引用记号 $x$，它一般是代表复数的。那么就要问：如果一个函数 $f(x)$ 自变量是复数 $z$，则
\[
f'(z_0)=\lim_{\Delta z\to0}\frac{f(z_0+\Delta z)-f(z_0)}{\Delta z}
\]
对不对？应该说对，只要右方的极限存在即可，但是复的 $\Delta z\to0$ 比之实的 $\Delta t\to0$ 内容要丰富得多。如果 $\Delta t$ 是实的，则 $\Delta t\to0$ 无非是从左侧或右侧趋于0，而当 $\Delta t$ 是复数时，$\Delta t\to0$ 的方式就是复杂得很难想象了。因此当 $z$ 为复数时，$\Delta f/\Delta z$ 有极限是一个极强的条件，所以我们就要问，$f(z)=e^z$ 有没有导数呢？复自变量函数的求导是一个很深刻的问题，下一章 §4 专门讨论它，而且恰好 $de^z/dz=e^z$ 总是对的。于是我们来看(23)式，并设其中 $t$ 是复的。如果 $t$ 之虚部为0，则 $t\to0$ 就意味着实变量 $t\to0$，这时(23)式就成为
\[
\lim_{t\to0}\frac{e^t-1}{t}=1.
\]
如果 $t$ 是纯虚的：$t=\sqrt{-1}x$，而实变量 $x\to0$，则因
\[
\frac{e^{\sqrt{-1}x}-1}{\sqrt{-1}x}=\frac{\cos x-1}{\sqrt{-1}x}+\frac{\sin x}{x},
\]
所以(23)式分开实虚部就给出两个极限式
\[
\lim_{x\to0}\frac{\cos x-1}{x}=0,
\qquad
\lim_{x\to0}\frac{\sin x}{x}=1.
\]
于是我们看到，微积分中两个基本极限其实是同一个极限的两个特殊情况。

但是这算不得一个证明，只是一条“道路”(approach)，要把它变成证明，还得补上复自变量函数求导的许多知识，说它是“道路”，因为它引导我们面对这许多新问题。

\textbf{3. 反三角函数}\quad 现在我们已经完成了通过微分方程来定义 $C(t)$ 与 $S(t)$ 并研究它们的性质，完全证实了它们就是中学学过的 $\cos t,\sin t$。所以下面我们不再使用 $C(t),S(t)$ 这两个记号了。同样 $\sqrt{-1}$ 也采用原来熟悉的 $i$ 来表示。

现在我们要讨论正、余弦函数的反函数问题，我们已知
\[
\frac{d\sin\theta}{d\theta}=\cos\theta. \tag{24}
\]
而当 $-\pi/2\le\theta\le\pi/2$ 时，由 $\cos\theta\ge0$（这一点由 $\cos\theta$ 是 $\overrightarrow{OP}$ 在 $x$ 轴上的投影即可看到）知 $\sin\theta$ 是 $\theta$ 的上升函数。它把 $[-\pi/2,\pi/2]$ 映到 $[-1,1]$。因此其反函数 $\theta=\theta(x)$ 存在，而且仍然是一个上升的 $C^1$ 函数（即一阶导数为连续的函数），映 $[-1,1]$ 到 $[-\pi/2,\pi/2]$。这样，读者当然会问，$\sin\theta$ 在 $[-\pi/2,\pi/2]$ 以外有没有反函数？当然也可以用类似上面的方法来处理，而例如又可得到 $[-1,1]\to[(k-1/2)\pi,(k+1/2)\pi]$，$k=0,\pm1,\pm2,\cdots$ 的 $C^1$ 函数，它们都应该是 $x=\sin\theta$ 的反函数 $\theta_k=\theta_k(x)$，也都定义在 $[-1,1]$ 上。但是其值一般不在 $[-\pi/2,\pi/2]$ 中，只有 $k=0$ 的一个例外：
\[
\widetilde\theta_0(x)=\theta(x).
\]
如果我们说有无穷多个反函数，则称 $k=0$ 的一个 $\widetilde\theta_0(x)=\theta(x)$ 为其主值，然后可以讨论这许多反函数之间有什么关系？这个关系式其实很简单，它就是
\[
\widetilde\theta_k(x)=k\pi+(-1)^k\widetilde\theta_0(x).
\]
它相应于正弦函数的一个基本关系式
\[
\sin[k\pi+(-1)^k\theta]=\sin\theta.
\]
但是我们不来讨论这个问题，我们也可以说 $\sin\theta$ 的反函数是一个“多值函数”，至于什么叫“多值函数”，下一节再谈。现在我们就只讲主值 $\widetilde\theta_0(x)$，我们习惯地记它为 $\arcsin x$ 而其它一切 $\theta_k(x)$ 则记为 $\mathrm{Arcsin}\,x$。

由(24)即有
\[
\frac{d\theta}{dx}=\frac1{\cos\theta}=\frac1{\sqrt{1-x^2}}. \tag{25}
\]
这里根号限取正值，因为对于主值 $\theta=\arcsin x$，$\cos\theta\ge0$ 且当 $x=0$ 时 $\theta=0$，所以
\[
\theta=\arcsin x=\int_0^x\frac{dx}{\sqrt{1-x^2}}. \tag{26}
\]
按道理说，$x$ 作为 $\sin\theta$ 之值就是 $[-1,1]$ 中的数，而 $\theta$ 作为一个角应该以弧度为单位，弧度可以用弧长来定义，但是更好是用扇形面积来定义。如图2-2-5，$\theta$ 之弧度值 $=2$ 扇形 $AOP$ 之面积，换一个记号，我们把(26)中的 $x$ 写成为 $y$，$y=\sin\theta$。为什么它称为“正弦”？它就是正对着角 $\theta$ 的半弦 $PR$，但是(26)是否扇形面积？这不要紧，我们来看积分 $\int_0^y\sqrt{1-y^2}\,dy$，它是曲边梯形 $OAPQ$ 的面积，用分部积分法即有
\[
\begin{aligned}
\int_0^y\sqrt{1-y^2}\,dy
&=y\sqrt{1-y^2}\big|_0^y+\int_0^y\frac{y^2\,dy}{\sqrt{1-y^2}}\\
&=y\sqrt{1-y^2}+\int_0^y\frac{(y^2-1)+1}{\sqrt{1-y^2}}\,dy\\
&=y\sqrt{1-y^2}-\int_0^y\sqrt{1-y^2}\,dy+\int_0^y\frac{dy}{\sqrt{1-y^2}}.
\end{aligned}
\]
移项后即有
\[
\int_0^y\sqrt{1-y^2}\,dy
=\frac12y\sqrt{1-y^2}+\frac12\int_0^y\frac{dy}{\sqrt{1-y^2}}.
\]
左边是 $OAPQ$ 面积，右方第一项是 $\triangle OPQ$ 之面积，所以
\[
\int_0^y\frac{dy}{\sqrt{1-y^2}}=2\text{ 扇形 }AOP\text{ 面积}=\theta\text{ 之弧度值}.
\]
这样一来，$y=\sin\theta$ 是由 $\theta$ 之弧度值算出正弦值，而 $\theta=\arcsin y=\int_0^y\frac{dy}{\sqrt{1-y^2}}$ 是由正弦之长（实际上是半弦）$y$ 算出其所对之角 $\theta$ 的弧度（按扇形面积理解）值。

\begin{figure}[htbp]
\centering
\begin{tikzpicture}[bella figure,scale=0.88,>=stealth]
  \draw[->] (-0.2,0)--(2.65,0) node[right] {$x$};
  \draw[->] (0,-0.2)--(0,2.65) node[above] {$y$};
  \draw (0,0) circle (1.8);
  \coordinate (P) at ({1.8*cos(40)},{1.8*sin(40)});
  \coordinate (R) at ({1.8*cos(40)},0);
  \coordinate (Q) at (0,{1.8*sin(40)});
  \draw (0,0)--(P)--(R); \draw[dashed] (P)--(Q);
  \draw (0.55,0) arc[start angle=0,end angle=40,radius=0.55];
  \node[right] at (P) {$P$}; \node[below] at (R) {$R$}; \node[left] at (Q) {$Q$}; \node[below left] at (0,0) {$O$};
  \node at (0.72,0.25) {$\theta$}; \node at (1.7,0.55) {$\sqrt{1-y^2}$}; \node at (0.18,1.0) {$y$}; \node at (0.95,1.0) {$1$};
  \node at (1.4,-1.15) {$x^2+y^2=1$};
\end{tikzpicture}
\caption*{图2-2-5}
\end{figure}

用这样的观点看自然对数 $\eta=\ln\xi$ 更有意义。上节我们说过，麦卡脱就发现了自然对数
\[
\eta=\int_1^\xi\frac{d\xi}{\xi}
\]
与双曲线的弧下的面积有关。由图2-2-6可见
\[
\int_1^\xi\frac{d\xi}{\xi}=\text{曲边梯形 }ABPQ\text{ 之面积}.
\]
但是
\[
\triangle OAQ=\frac12\cdot1\cdot1=\triangle OBP=\frac12\xi\cdot\eta=\frac12,
\]
双方同减去 $\triangle OAT$，即有 $\triangle OTQ=\text{梯形 }ABPT$，双方再加上 $TPQ$ 即有
\[
\text{扇形 }OPQ\text{ 面积}=\text{曲边梯形 }ABPQ\text{ 之面积}.
\]
所以，对数函数 $\ln\xi=\int_1^\xi d\xi/\xi$ 作为面积与 $\arcsin x$ 是十分相近的，如果我们作一个旋转
\[
\xi=\frac1{\sqrt2}x+\frac1{\sqrt2}y,
\qquad
\eta=\frac1{\sqrt2}x-\frac1{\sqrt2}y,
\]
则 $\xi\eta=1$ 变成 $x^2-y^2=2$，则这种相似更为明显。不过，对于反三角函数
\[
\text{扇形 }AOP\text{ 面积}=\frac12\int_0^y\frac{dy}{\sqrt{1-y^2}},
\]
式中多了一个因子 $1/2$，是由于图2-2-5中的圆是 $x^2+y^2=1$，而图2-2-6中的双曲线是 $x^2-y^2=2$。如果注意到若在圆的方程中把 $y$ 换成 $iy$，则圆就变成双曲线，所以一旦进入复域，正弦—反正弦就会变成指数—对数。欧拉公式也就不足为奇了。

\begin{figure}[htbp]
\centering
\begin{tikzpicture}[bella figure,scale=0.82,>=stealth]
  \draw[->] (0,0)--(4.9,0) node[right] {$\xi$};
  \draw[->] (0,0)--(0,4.6) node[above] {$\eta$};
  \draw[domain=0.45:4.4,smooth,variable=\x,thick] plot ({\x},{1.7/\x});
  \coordinate (Q) at (1.3,{1.7/1.3});
  \coordinate (P) at (3.4,{1.7/3.4});
  \draw[dashed] (Q)--(1.3,0) node[below] {$A$};
  \draw[dashed] (P)--(3.4,0) node[below] {$B$};
  \draw (0,0)--(Q); \draw (0,0)--(P);
  \draw[->] (0,0)--(3.7,3.7) node[right] {$x$};
  \draw[->] (0,0)--(-1.8,1.8) node[left] {$y$};
  \node[above right] at (Q) {$Q(1,1)$}; \node[right] at (P) {$P(\xi,\eta)$};
  \node[below left] at (0,0) {$O$};
\end{tikzpicture}
\caption*{图2-2-6}
\end{figure}

当年麦卡托正是看到了这一点，然后传到牛顿才得到许多初等函数的幂级数展开式。欧拉在他的《无穷小分析引论》一书中也就把以 $e$ 为底的对数称为“自然对数或双曲对数”。

\textbf{4. 振动现象}\quad 以上我们从匀速圆周运动向两个互相垂直的坐标轴投影，找到了 $z(t),C(t)$，以及 $S(t)$ 所适合的常微分方程。它们的相互关系给出了十分重要的欧拉公式。$z(t),C(t)$，和 $S(t)$ 都是适合同样类型的微分方程
\[
\ddot x+\alpha^2x=0 \tag{27}
\]
（但初始条件不同）的函数。由于它在物理和工程中的重要性，我们称适合(27)的 $x(t)$ 为一个谐振子(harmonic oscillator)，所以 $z(t)$ 是一个复的谐振子，而 $C(t),S(t)$ 则是实谐振子。实际用起来，则复谐振子更方便。

一个一般的复谐振子总是 $e^{i\alpha t}$ 与 $e^{-i\alpha t}$ 的复系数的线性组合，但是只要弄清楚了其中之一，则另一个自然完全都清楚了。所以下面我们限于讨论 $Ce^{i\alpha t}$，这里 $C$ 是一个复数，但是由欧拉公式我们很容易看到任一复数 $C$ 都可写为 $Ae^{i\varphi}$，这里 $A>0$，所以谐振子一定可以写为
\[
Ae^{i(\alpha t+\varphi)}=Ae^{i(2\pi\omega t+\varphi)}. \tag{28}
\]
这里 $A$ 称为其振幅，$\alpha$ 称为圆频率，$\omega$ 称为频率，$\omega$ 之单位就是我们很熟悉的赫兹(Hertz)，每秒振动一次称为1 Hz，量纲则是 $1/T$，所以 $\alpha t$ 或 $\omega t$ 都是无量纲量\jiaokan{原文此处再次将 $\alpha$、$\omega$ 的“频率/圆频率”名称及赫兹单位对调；按文中 $\alpha=2\pi\omega$ 的定义改正。}。$\alpha t+\varphi=2\pi\omega t+\varphi$ 称为相位(phase)，这样一来，什么叫相位有了一个数学的说法：今后我们时常会遇到形如 $A(x,t)e^{i\Phi(x,t)}$ 的量，而且 $A(x,t)\ge0$，这时我们总是说 $A(x,t)$ 是振幅，$\Phi(x,t)$ 是相位，$\Phi(x,t)|_{t=0}$ 称为初（始）相（位），所以对于(28)，$\varphi$ 就是初始相位。

谐振子为什么如此重要？在自然界中我们常见到这样的运动：某一物体（某一质点，某一“电”量，如电位、电容、……）在某一位置附近往复振动。其例子不仅有摆，还有如电子在原子中的振动，所谓脉冲星一阵一阵地例如以 X 射线形式发出大量的能量；激光在激光器内往复反射……，还有电路中的电流、电压……心电图、血压……这些都是振动，然后这些振动在空间中的传播，就成为波。例如水的振动在水面上传播成为水波，天线中的电磁振动在空间传播为电磁波，电压、电流的振动在导线中传播输送出电讯号。不但自然界如此，社会生活中也有越来越多的现象可以归入其中而用数学方法考察。例如价格、股市、流行病的发生等等都是。尽管研究的对象可以几乎是包罗万象，研究它们的数学方法却是统一而有系统的。常常总是把它们归为一个微分方程（但时常又是一个差分方程），而(27)可以说是最简单的一例。

(27)总是一个理想的情况，称为一维谐振子。说它是一维的，因为刻画这个物理现象只需一个物理量 $x(t)$（时间 $t$ 总是作为参数来看待的），物理学中许多重要的理论模型都是建筑在它的基础上的。例如考虑单摆（图2-2-7），就是质量 $m$ 的质点挂在长为 $l$ 的弦上并在重力下振动。要想刻画它只需要知道一个物理量：随 $t$ 的变化即可，所以它是一维振动。它是理想的，因为我们假设弦的质量可以忽略不计，而且是刚性的，因此 $l$ 之大小不变。这时质点只受重力 $mg$（向下）作用，重力有两个分力，一是在弦的方向上，因为弦是刚性的，这个分力不起作用。二是在圆上的圆弧的切线方向上，大小为 $mg\sin x$ 而且指向平衡位置，所以必与角位移 $x(t)$ 方向相同，但符号相反，称为恢复力。我们考虑单摆的小振动，就是角位移之绝对值 $|x(t)|$ 很小，所以 $\sin x\sim x$，这样恢复力成为 $-mgx$。牛顿的第二定律给出
\[
ml\ddot x(t)=-mgx(t).
\]
左方出现 $l$ 是因为 $lx$ 为线位移，而第二定律的加速度并非角加速度 $\ddot x(t)$，而是线加速度 $l\ddot x(t)$。此即方程(27)。

\begin{figure}[htbp]
\centering
\begin{tikzpicture}[bella figure,scale=0.85,>=stealth]
  \coordinate (O) at (0,3.6); \coordinate (B) at (0,0.55); \coordinate (M) at (1.35,0.9);
  \draw[thick] (O)--(B); \draw[thick] (O)--(M) node[midway,right] {$l$};
  \draw[dashed] (B) arc[start angle=-90,end angle=-66,radius=3.05];
  \draw (O) ++(-90:0.75) arc[start angle=-90,end angle=-66,radius=0.75];
  \node at (0.27,2.72) {$x$}; \node[above] at (O) {$O$}; \fill (M) circle (2.2pt); \node[right] at (M) {$m$};
  \draw[->] (M)--++(0,-1.2) node[below] {$mg$};
  \draw[->] (M)--++(-0.58,-0.28) node[below left] {$mg\sin x$};
\end{tikzpicture}
\caption*{图2-2-7}
\end{figure}

上面我们研究三角函数时，是把 $\cos t=C(t)$ 放在 $x$ 轴上，$\sin t=S(t)$ 放在 $y$ 轴上，再考虑一个复平面。$x$ 轴与 $y$ 轴是圆上的点的投影运动的场所，我们不妨称之为物理空间。在单摆的例子中，我们也不妨把角位移 $x(t)$ 放在 $x$ 轴上，这个 $x$ 轴就应称为物理空间（一个恰当的数学名词是构形空间(configuration space)，因为它的一个点 $x$ 就给出了角的大小，因而决定了摆的形状——即构形），那么是什么相应于 $y$ 轴呢？回到三角函数的情况，$x$ 轴既用来表示 $\cos t$，$y$ 轴则表示 $\sin t=-d\cos t/dt$。现在在单摆的情况下是否也可以用 $y$ 轴来表示 $\dot x(t)$ 即角速度呢？（这里相差一个符号，但它显然没有影响。）或者用来表示 $p=m\dot x$ 即动量呢？当然都是可以的，真正重要的是要看动量 $p$ 与位移 $x$ 是怎样联系起来的。由方程(27)——其中 $\alpha^2=g/l$——双方乘以 $\dot x$ 再积分，有
\[
\begin{aligned}
0&=\int_0^t(\dot x\ddot x+\alpha^2\dot x x)\,dt
=\frac12\int_0^t\frac d{dt}(\dot x^2+\alpha^2x^2)\,dt\\
&=\frac12(\dot x^2(t)+\alpha^2x^2(t))-\frac12(\dot x^2(0)+\alpha^2x^2(0)).
\end{aligned}
\]
所以 $\dot x^2(t)+\alpha^2x^2(t)$ 是守恒的。但是我们要注意，除了相差一个常数因子（其中有质量 $m$，还有 $g/l$），恢复力是
\[
-\frac d{dx}U(x)=-\frac d{dx}\left(\frac{\alpha^2x^2}{2}\right),
\]
所以 $\alpha^2x^2/2$ 乃是质点的位能，另一方面 $m\dot x^2/2$ 是动能，而
\[
\frac{p^2}{m^2}+\alpha^2x^2=\dot x^2(t)+\alpha^2x^2(t)=\text{常数} \tag{29}
\]
就表示能量守恒：
\[
\text{(动能)}+\text{(位能)}=\text{常数}.
\]
现在我们就作一个平面：$x$ 轴表示位移 $x(t)$，$y$ 轴表示动量 $p(t)=m\dot x(t)$，并称此平面为相平面（相空间）(phase plane, phase space) 而与构形空间相区别；$(x,p)$ 就称为此质点的相，因为有此就足以完全决定此质点的运动状况。$x(t)$ 的更高阶的导数不必给了，因为牛顿运动方程就给出了 $\ddot x(t)$，而更高阶的导数也可以从运动方程求导得到。不但现时该质点的运动状况完全决定了，即其过去未来的运动状况也可以用此时的 $(x,p)$ 作为初始条件再通过方程(27)完全决定。当 $x(t),p(t)$ 已经决定了以后，$(x(t),p(t))$ 就给出了相空间的轨道，称为相轨道，而与构形空间中的轨道 $x=x(t)$ 相区别。于是对于单摆的小振动，相轨道就是椭圆(29)，亦即\jiaokan{原文下式第二项印作 $\alpha^2x(t)$，与式(29)及“相轨道就是椭圆”的叙述矛盾；应为 $\alpha^2x^2(t)$。}
\[
\frac{p^2(t)}{m^2}+\alpha^2x^2(t)=C.
\]
当 $m=\alpha=1$ 时就是圆。于是，从力学的观点看来，这里的相平面就是前面讲的 $z$ 平面的一般化，而前面讲的
\[
\cos^2t+\sin^2t=1,
\]
从几何上看无非是勾股定理，而从力学角度来看就是能量守恒定律的特例。真正值得注意的是，这里得出了椭圆而前面得出了圆，差别仅来自恢复力系数 $\alpha^2$ 之大小。从拓扑学角度来看，它们并无本质的区别，都是封闭曲线，所以都表示周期运动。下一章我们要详细讨论开普勒第三定律与牛顿万有引力定律的关系，到那时我们会再一次看到行星的运动也包含在这样的框架内。

现在再进一步设质点在阻力亦即摩擦力下运动。阻力是一个复杂的对象，在最简单的情况下，可以假设阻力为 $-h\dot x$，$h>0$ 是一个常数，这里加了一个负号是因为阻力总与速度方向相反，否则就谈不上阻滞了。这时由牛顿第二定律给出
\[
\ddot x+h\dot x+\alpha^2x=0. \tag{30}
\]
（我们略去了质量 $m$，而把系数作必要的修正）。我们首先把(30)化成一个方程组。为此令 $x_1=x$，再引入 $x_2=\dot x_1$，则(30)化成
\[
\dot x_1=x_2,
\]
\[
\dot x_2=-hx_2-\alpha^2x_1. \tag{31}
\]
$(x_1,x_2)$ 平面现在就是相平面。所以，研究方程组(31)即是在相平面上讨论它；但是为了求解，我们仍用(30)，并且仍设 $x(t)=Ce^{\lambda t}$，代入(30)，即得关于 $\lambda$ 的一个二次方程
\[
\lambda^2+h\lambda+\alpha^2=0 \tag{32}
\]
称为特征方程，它有两个特征根
\[
\lambda_{1,2}=-\frac h2\pm\frac12\sqrt{h^2-4\alpha^2}.
\]
这时最重要的是 $h^2-4\alpha^2$ 的符号。下面我们只看 $\lambda_1$，因为 $\lambda_2$ 是 $\lambda_1$ 的复共轭，它不会给出新的物理知识。

相应于 $\lambda_1$，将得到解
\[
x(t)=Ce^{-\frac h2t}e^{i\widetilde\alpha t}, \tag{33}
\]
\[
\widetilde\alpha^2=\alpha^2-h^2/4.
\]
$h=0$，就是没有阻力的情况，如果 $h$ 逐步增大，把 $C$ 写成一个复数 $C=Ae^{i\varphi_0},A>0$，有
\[
x(t)=Ae^{-\frac h2t}e^{i(\widetilde\alpha t+\varphi_0)}.
\]
于是我们看到频率由 $\alpha$ 变成 $\widetilde\alpha$，就是频率变慢，振幅出现了因子 $e^{-ht/2}$，而当 $t\to\infty$ 时衰减为0。但是它仍然表现了振动性质。$h=2\alpha$ 是一个临界值，这时 $\widetilde\alpha=0$，但是当特征方程出现重根时，解不再写成(33)而有变化。如果阻力再增加，使 $h^2-4\alpha^2>0$，解 $\lambda_1$ 与 $\lambda_2$ 都成了实数，而解(33)现在成了
\[
x(t)=Ae^{\lambda_1t+i\varphi_0},
\qquad
\lambda_{1,2}=-\frac h2\pm\frac12\sqrt{h^2-4\alpha^2}<0.
\]
这时就不再有振动，而只有指数衰减了。

以上我们是求出了 $x(t)$，如果转到相空间 $(x_1,x_2)$，尽管 $x_1(t),x_2(t)$ 都不难得出，却不容易找到与椭圆(29)相应的联结 $x_1,x_2$ 的关系式。实际上，相轨道不会再是一个封闭曲线，因为不再有周期运动。但是相轨道的分析引导到数学与物理学中、工程学中十分富有成果的分支。
